\begin{musica}{Guantanamera}{José Martí, Josito Fernandez}{Tuna Universitária do Porto}
\nota{Instrumental: Dmaj7 E A Dmaj7 E}

\nota{Refrão:}
\ac{A}{}_ Guanta\ac{D}{na}me\ac{E}{ra}, \ac{D}{}_ guajira \ac{A}{guan}tana\ac{D}{me}ra, \ac{E D}{}_ _
\ac{A}{Guan}tana\ac{Bm7}{me}r\ac{E D}{a}, _ guajira \ac{A}{guan}tana\ac{D}{me}r\ac{E}{a}.

\ac{D}{Yo} soy un \ac{A}{hom}bre \ac{D}{sin}ce\ac{E}{ro}
\ac{D}{De} donde \ac{A}{cre}c\ac{D}{en} las \ac{E}{pal}mas.
\ac{D}{Yo} soy un \ac{A}{hom}bre \ac{D}{sin}ce\ac{E}{ro}
\ac{D}{De} donde \ac{A}{cre}c\ac{D}{en} las \ac{E}{pal}mas.
\ac{D}{Y} antes \ac{A}{de} mo\ac{D}{rir}me \ac{E}{qui}ero
\ac{D}{E}char mis \ac{A}{ver}sos del \ac{D}{al}m\ac{E}{a}. \ac{D}{}_

\nota{Refrão}

Mi verso es de un verde claro
Y de un carmin encendido,
Mi verso esde un verde claro
Y de un carmin encendido.
\ac{Amaj7}{Mi} verso es un \ac{Bm}{cier}vo he\ac{C#m7 Cm7}{ri}do ___
\ac{Bm7}{Que} busca en el monte ampa\ac{E}{ro}.

\nota{Refrão}

\ac{Bm7}{Guan}taname\ac{E}{ra}, \ac{E7}{}_ guajira \ac{A}{guan}tana\ac{Bm7}{me}ra, \ac{E E7}{}_ _
\ac{A}{Guan}tana\ac{Bm7}{me}r\ac{E}{a}, \ac{E7}{}_ guajira \ac{F#}{guan}tana\ac{B7}{me}r\ac{E}{a}.

\ac{A}{Cu}ltivo \ac{Bm7}{una} rosa \ac{E}{bla}n\ac{E7}{ca}
\ac{A}{En} Ju\ac{Bm7}{nio} com en \ac{E}{E}ne\ac{E7}{ro}
\ac{A}{Pa}ra el amigo \ac{Bm7}{sin}ce\ac{E}{ro} \ac{E7}{}__
\ac{A}{Que} me da su \ac{Bm7}{ma}no \ac{E}{fra}n\ac{E7}{ca}.
Para el a\ac{A}{mi}go \ac{Bm7}{sin}ce\ac{C#m7}{ro}
\ac{Cm7}{Que} me da \ac{Bm7}{su} mano fran\ac{E}{ca}.

\posicoes{Cm: x35543, C#m7: x35343, D: xx0232, E: 022100, F#: 244322, A: x02220,
  Amaj7: x02120, B7: 021202, Bm: x24432, Bm7: x24232}
\end{musica}
